\subsection{The complete cubic test for one pointed quartic in characteristic five}

\begin{claim}{B.80}{Proved}{All cubic covers of the pointed quartic $W^4=X^3Z-XZ^3$}
Let
\[
 C/\mathbf F _5:\quad W^4=X^3Z-XZ^3,
 \qquad P=(1:0:0),\qquad K=\mathbf F _{5^8}.
\]
Every nonzero geometric order-three character first becomes constant
over $K$. There are $364$ connected constant cyclic degree-three
$P$-split covers of $C_K$, up to inversion of the deck character.
Exactly $96$ are class-number witnesses; the other $268$ fail the
test over $K$ and every further constant extension. The successful
ones remain witnesses over $\mathbf F _{5^{8n}}$ exactly when
$3\nmid n$. Consequently the least admissible odd prime giving a
witness for this pointed quartic over some extension is three, and
the least field for an order-three witness is $\mathbf F _{5^8}$.

More explicitly, the three elliptic quotients identify
$J_C[3]$ over $\overline{\mathbf F}_5$ with
$E_0[3]\oplus E_1[3]\oplus E_2[3]$. A cover represented by
$(T_0,T_1,T_2)$ is successful precisely when exactly two entries
are nonzero. On its new coherent cohomology, the characteristic
polynomial of the $K$-linear Cartier iterate is
\[
 \begin{cases}
 (T^2+T+1)^2,&\text{exactly two nonzero entries},\\
 (T-1)^4,&\text{one or three nonzero entries}.
 \end{cases}                                                   \tag{B.80.1}
\]
The different optimization in B.34, minimizing the extension degree
first among all primes, has degree three and least prime thirteen
there. It is not the prime-first minimum asserted here.
\end{claim}

\begin{proof}
We first justify that the calculation covers every geometric character.
Choose $\eta$ with $\eta^4=-1$, so $\eta\in\mathbf F _{25}$.
The change of variables
\[
 U=X-2Z,\qquad V=W,\qquad T=\eta(X+2Z)
\]
takes $C$ to the Fermat quartic $U^4+V^4+T^4=0$.
Its three sign involutions generate a Klein four group, whose common
quotient is the genus-zero conic in $U^2,V^2,T^2$.
The three individual quotients are elliptic curves. Their quotient
maps $q_i:C\to E_i$ give homomorphisms
\[
 \phi:\prod_iE_i\to J_C,\quad (a_i)\mapsto\sum_iq_i^*a_i,
 \qquad
 \psi:J_C\to\prod_iE_i,\quad a\mapsto(q_{i*}a)_i.
\]
On $J_C$, $q_i^*q_{i*}=1+\iota_i$. The norm of the Klein four
group factors through the genus-zero quotient, and is zero. Hence
$\phi\psi=\sum_i(1+\iota_i)=2$. The mixed maps
$q_{i*}q_j^*$ vanish for $i\ne j$: the involution $\iota_i$ acts
by $-1$ on the elliptic quotient attached to $\iota_j$, since its
quotient there is the genus-zero common quotient. The mixed map is
therefore killed by two and must vanish as a homomorphism of abelian
varieties. Thus $\psi\phi=2$ as well. In particular the maps induce
inverse isomorphisms, up to multiplication by two, on three-torsion.

All three elliptic quotients become isomorphic over $\mathbf F _{25}$
to $E:y^2=x^3-x$. This follows explicitly by permuting the Fermat
coordinates and reversing the displayed linear change. The three
matrices implementing it are retained, and substitution checks each
quartic equation identically. The curve $E$ has eight points over
$\mathbf F _5$, so its Frobenius satisfies
$F_E^2+2F_E+5=0$. Modulo three this gives $F_E^4=-1$ and
$F_E^8=1$. Since the isogeny above is defined over $\mathbf F _{25}$,
the same identities hold on $J_C[3]$. No nonzero vector can have
period dividing four. Every period divides eight, so every nonzero
vector has period exactly eight. On the character space the dual
contragredient action has the same property. Thus degree eight is
the first field on which any nonzero order-three character is constant.

We describe all of the covers and their pointings. For $T_i=(a_i,b_i)$
in the nonzero three-torsion of $E_i$, let
\[
 m_i=\frac{3a_i^2-1}{2b_i},\qquad c_i=b_i-m_i a_i,
 \qquad g_i=y_i-m_ix_i-c_i.
\]
The tangent is a flex, so
$\operatorname{div}(g_i)=3T_i-3O_i$. Directly,
$(m_i(x-a_i)+b_i)^2-(x^3-x)=-(x-a_i)^3$.
For an active set of entries put $G=\prod_iq_i^*g_i$ and divide
it by its leading unit at $P$. The normalization is explicit in the
certificate. If the identity quotient is active, it contributes a
pole of order six at $P$, with leading coefficient one; this pole
is a cube and gives no ramification. All other active factors are
units at $P$. The normalization of
\[
                             z^3=G
\]
is therefore an \emph{\'etale}, constant cyclic, $P$-split cover.
It is connected because the corresponding nonzero three-torsion
line bundle is not trivial. Replacing every $T_i$ by $-T_i$ reverses
the deck generator and gives the same cover. These are the only
identifications. Consequently the $9^3-1$ nonzero triples give
$(9^3-1)/2=364$ covers, with respectively
\[
               12,\quad 96,\quad 256
\]
having one, two, or three nonzero entries. The equality with all
constant characters uses $\mu_3\subset K$ and the Kummer
identification with $J_C[3]$; the chosen unit normalization removes
the constant twist ambiguity.

We now justify the differential calculation. In the affine coordinate
ring $K[x,w]/(w^4-x^3+x)$, write
$q_i=(X_i/Z_i,W_i^2/Z_i^2)$, with the displayed linear forms coming
from the three coordinate matrices. Set
$Q_i=W_i^2-m_iX_iZ_i-c_iZ_i^2$. If $s$ entries are active, take
homogeneous forms $H$ of degree $s+1$. In the two nontrivial deck
characters, regular differentials have the form
\[
                \frac{H\,dx}{w^3\prod_iZ_i\,z^a},
                         \qquad a=1,2.                       \tag{B.80.2}
\]
To see the precise regularity conditions, let
$D_i=q_i^{-1}(T_i)$ and $H_i=q_i^{-1}(O_i)$. The line $Z_i=0$
cuts the divisor $2H_i$ on $C$, and $\operatorname{div}(G)$ is
$3\sum_i(D_i-H_i)$. The differential $dx/w^3$ represents the
canonical plane section. Therefore (B.80.2) is regular exactly when
the section $H$ vanishes along
\[
                         \sum_i(aD_i+(2-a)H_i).
\]
For $a=1$ these are simple conditions at $D_i$ and double conditions
at the hyperflex supporting $H_i$; for $a=2$ they are double conditions
at $D_i$. All these conditions are explicit linear equations in the
coefficients of $H$. Each solution space has dimension two:
Riemann--Roch gives $g(C)-1=2$ for the canonical bundle tensored
with a nontrivial degree-zero three-torsion line bundle.
The programs construct two independent forms, verify every vanishing
and derivative condition, and hence obtain a full basis for each
character. Together they give the new four-dimensional differential
space; the genus of each cover is seven.

The first program computes Cartier by extracting the coefficients
numbered $5j+4$ after clearing denominators and taking fifth roots.
Every division and basis-coordinate identity is checked with zero
remainder. A separate program imports no producer code and checks
the different rational-function identity
\[
 \left(\frac{\mathcal C(v\,dx)}{dx}\right)^5
                          =-\left(\frac d{dx}\right)^4v.
                                                               \tag{B.80.3}
\]
Identity (B.80.3) follows by writing
$v=\sum_{j=0}^4v_j^5x^j$: the fourth derivative kills every term
except $4!v_4^5=-v_4^5$. Thus it determines Cartier exactly.
In our function field $dw/dx=(3x^2-1)/(4w^3)$; the verifier applies
this derivation to every component in the basis $1,w,w^2,w^3$.
All $1456$ Cartier columns and all $5824$ component identities pass.
It also checks the eightfold semilinear product in the opposite
implementation order and calculates the characteristic polynomial
by a separate recursive adjugate formula. The result is (B.80.1) for
every one of the $364$ covers. The exact field modulus, torsion
points, coordinate matrices, cover normalizations, differential
bases, Cartier matrices and their eighth iterates are all retained.

By duality, the linear Cartier iterate and coherent Frobenius have
the same characteristic polynomial. The Artin--Schreier sequence
identifies the Frobenius-stable coherent part with the geometric
five-torsion cohomology, while its nilpotent part contributes no
factor to the determinant at one. Since the cover degree is prime
to five, pullback and norm split the old and new five-primary data.
Thus the relative class number is a unit exactly when the polynomial
in (B.80.1) is nonzero at one. This is true precisely in the two-entry
case. Under a further constant extension of degree $n$, the two
primitive cube-root eigenvalues acquire an eigenvalue one exactly
when $3\mid n$. A pre-existing eigenvalue one in either of the
other cases stays one under every power. This proves all extension
claims, and three is the smallest admissible odd prime in
characteristic five.
\end{proof}

\paragraph{Check record.}
Checked only in this writing as a promoted statement: GPT-6 Astra Pro,
finite-checks reader, 18 November 2026. We read both complete program
sources and verified the mathematical interpretation above, reran the
producer and its independent fourth-derivative verifier, and compared
all $364$ results. These are reruns of the supplied programs, not a
claim that their code was newly independently reimplemented here.
The effort setting is not separately exposed by the runtime.

\paragraph{Remaining scope of B6.}
This theorem concerns exactly the displayed pointed quartic. It does
not classify every quartic over $\mathbf F _3$ or $\mathbf F _5$,
and does not compare the unrestricted existence of abelian and
nonabelian witnesses for arbitrary curves. Those quantifiers are
retained in B.32.
